\scanpage{188}
\mspage{133}
% archive sheet number: 187

pour \deleted{tout} $a$ dans $\mathcal A=\mathbb D(e)$. (\deleted{\unsure{D'où}} $p:S\longrightarrow e$, $p':S'\longrightarrow e$ sont les morph. structuraux sur la cat. ponctuelle.) On a vu la formule heuristique générale

\[
p_!p^*\underbrace{(a)}_{a_S}\unsure{(\!=\,\prime|=}\,\mathrm{hot}(S)\times a,
\]

où $\mathrm{hot}\,S$ est la classe de $S\in\mathrm{Ob}\,\mathrm{Cat}$ dans $\mathrm{Hot}^{\unsure{?|S}}=(\mathrm{Cat})W_\infty^{-1}$, et où le produit $\times$ est une opération externe de $\mathrm{Hot}$ sur $\mathcal A=\mathbb D(e)$

\[
\mathrm{Hot}\times\mathcal A\longrightarrow\mathcal A.
\]

Dans le cas actuel, $\mathcal A=\mathrm{Hot}$, et l'opération

\[
\mathrm{Hot}\times\mathrm{Hot}\longrightarrow\mathrm{Hot}
\]

ne doit être autre que le \deleted{produit} foncteur ``produit cartésien'', \inserted{lequel est} déduit du foncteur similaire dans $\mathrm{Cat}$. Quand on dispose de ce formalisme, le résultat cherché pour un $f:S'\longrightarrow S$ se réduit simplement à dire que si $f$ est dans $W$, alors il en est de $=$ de $f\times\mathrm{id}_X$ pour tout $X$ dans $\mathrm{Cat}$ (en prenant $X$ tel que $\mathrm{hot}(X)\simeq a$). Or là c'est un fait bien connu, résultant des axiomes mis sur $W$.

J'ai donc l'impression que je n'ai aucune chance de prouver Der 3 a$'$ (\unsure{-p.|p.} 111), ni même Der 3 a) (cas où $S=e$), sans avoir une maîtrise du formalisme des $p_!$. Et j'ai
